% 01_introduction.tex

\section{Introduction}
\label{sec:intro}

No-reference mean opinion score (MOS) predictors are increasingly used to screen and compare speech systems at scale. DNSMOS P.835~\cite{Reddy2022dnsmos}, NISQA~\cite{Mittag2021nisqa}, and UTMOS~\cite{Saeki2022utmos} reduce reliance on repeated large-scale listening tests, while DNSMOS also serves as an official metric in the DNS Challenge~\cite{dubey2024icassp,Reddy2019scalable}. Their use has outpaced scrutiny of their failure modes.

Existing MOS predictors return point estimates without an explicit reliability measure. A practitioner therefore cannot tell from the score alone whether an utterance should be trusted, reviewed, or sent to a listening test. Feature-space distance is an attractive post-hoc signal because it requires neither predictor retraining nor labelled shifted data. Its interpretation is not automatic, however. A score that recognizes a change in language, channel, or degradation process may still fail to identify which individual MOS predictions are inaccurate.

We test this hypothesis through a conformal uncertainty wrapper for frozen MOS predictors. Split conformal prediction~\cite{Angelopoulos2023conformal} provides finite-sample marginal coverage when calibration and test examples are exchangeable and calibration labels are independent of model fitting. The wrapper first constructs a global interval from calibration residuals. It then applies an established Mondrian-style construction~\cite{Bostrom2020Mondrian} in which Mahalanobis distance~\cite{Lee2018} selects a distance-specific interval radius. The central question is whether a feature-space domain-shift signal transfers to utterance-level error ranking strongly enough to support this adaptation.

This work makes three contributions. First, we construct a leakage-controlled evaluation with disjoint feature-reference, conformal-calibration, and in-domain evaluation roles. Second, we compare Euclidean, Mahalanobis, and nearest-neighbour scores across the distinct tasks of domain detection and utterance-level error detection for three MOS predictors. Third, we test the practical objective directly by comparing fixed and distance-stratified conformal intervals under five metadata-defined shifts. Mean domain-detection AUROC is 0.615--0.690, while mean AUROC for detecting the highest-error quartile is only 0.433--0.499. Distance-stratified interval adaptation is predictor and domain dependent. The contribution is therefore both methodological and empirical, a controlled test showing why an OOD score must be validated as a signal of prediction error, not merely of domain membership, before it is used to adapt prediction uncertainty.
